\documentclass[11pt]{article}

% ── Packages ──────────────────────────────────────────────────────────
\usepackage[utf8]{inputenc}
\usepackage[T1]{fontenc}
\usepackage[margin=1in]{geometry}
\usepackage{parskip}
\usepackage{listings}
\usepackage{xcolor}
\usepackage{hyperref}
\usepackage{tcolorbox}
\usepackage{tabularx}
\usepackage{booktabs}
\usepackage{amsmath}
\usepackage{amssymb}
\usepackage{enumitem}
\usepackage{fancyvrb}
\usepackage[hang,flushmargin]{footmisc}
\usepackage[expansion=false]{microtype}
\DisableLigatures{encoding = T1, family = tt*}

% ── Hyperref setup ─────────────────────────────────────────────────────
\hypersetup{
    colorlinks=true,
    linkcolor=blue!50!black,
    urlcolor=blue!50!black,
    citecolor=blue!50!black,
    pdftitle={Single Liposome Curvature Assay Protocol},
}

% ── Listings setup (code blocks) ──────────────────────────────────────
\definecolor{codebg}{RGB}{248,248,248}
\definecolor{codeframe}{RGB}{220,220,220}

\lstdefinestyle{terminal}{
    backgroundcolor=\color{codebg},
    basicstyle=\ttfamily\small,
    breaklines=true,
    breakatwhitespace=false,
    columns=fullflexible,
    keepspaces=true,
    showspaces=false,
    showstringspaces=false,
    frame=single,
    framesep=4pt,
    rulecolor=\color{codeframe},
    xleftmargin=8pt,
    xrightmargin=8pt,
    aboveskip=8pt,
    belowskip=8pt,
}

\lstset{style=terminal}

% Inline code formatting
\newcommand{\code}[1]{\texttt{#1}}
\newcommand{\flag}[1]{\texttt{#1}}

% ── Callout boxes ─────────────────────────────────────────────────────
\newtcolorbox{notebox}{
    colback=blue!5,
    colframe=blue!40,
    boxrule=0.5pt,
    arc=2pt,
    left=8pt, right=8pt, top=6pt, bottom=6pt,
}

\newtcolorbox{warnbox}{
    colback=orange!5,
    colframe=orange!50,
    boxrule=0.5pt,
    arc=2pt,
    left=8pt, right=8pt, top=6pt, bottom=6pt,
}

\newtcolorbox{tipbox}{
    colback=green!5,
    colframe=green!50!black!50,
    boxrule=0.5pt,
    arc=2pt,
    left=8pt, right=8pt, top=6pt, bottom=6pt,
}

% ── Section formatting ────────────────────────────────────────────────
\usepackage{titlesec}
\titleformat{\section}{\Large\bfseries\color{blue!50!black}}{\thesection}{0.6em}{}
\titleformat{\subsection}{\large\bfseries\color{blue!30!black}}{\thesubsection}{0.6em}{}
\titleformat{\subsubsection}{\normalsize\bfseries}{\thesubsubsection}{0.5em}{}
\titlespacing*{\section}{0pt}{18pt}{6pt}
\titlespacing*{\subsection}{0pt}{14pt}{4pt}

% ── Title ──────────────────────────────────────────────────────────────
\title{
    \vspace{-1cm}
    \Huge \textbf{Single Liposome Curvature Assay} \\[0.3cm]
    \LARGE Image Analysis Protocol \\[0.2cm]
    \large \textit{From raw microscopy images to curvature-sorting plots}
}
\author{}
\date{Version 2.0 \\ For users with no prior coding experience}

% ──────────────────────────────────────────────────────────────────────
\begin{document}

\maketitle
\thispagestyle{empty}

\vspace{1cm}

\begin{notebox}
\textbf{Who this is for.} This document is for wet-lab users running
the SLiC analysis pipeline for the first time, with no prior Python or
command-line experience. Every command is shown for both Mac/Linux and
Windows PowerShell. If you are already comfortable with Python and the
terminal, the \texttt{README.md} file in the repository is a more
concise reference.
\end{notebox}

\newpage

\tableofcontents
\newpage

% ═══════════════════════════════════════════════════════════════════
\section{Overview}

This protocol takes you from raw two-channel fluorescence microscopy
images all the way to a plot of \textbf{protein surface density vs.\
liposome diameter} (curvature). It is the computational half of the
single-liposome curvature (SLiC) assay.

\subsection*{What the assay measures}

Each liposome in your image appears as a sub-diffraction-limit
fluorescent spot. The \textbf{lipid channel amplitude} is proportional
to the liposome's surface area, so its square root is proportional to
diameter. The \textbf{protein channel amplitude} tells you how much
protein is bound. By calibrating fluorescence amplitude against an
independent dynamic light scattering (DLS) measurement of liposome
size, you get a conversion factor from amplitude to physical diameter.
The protein surface density is then
\[
\rho_{\text{protein}} = \frac{A_{\text{protein}}}{\pi D^2}
\]
which you plot versus diameter. A curvature-sensing protein shows
$\rho_{\text{protein}}$ rising at small diameters (high curvature).

\subsection*{The four steps}

\begin{enumerate}[leftmargin=*]
    \item \textbf{Prepare images} --- split microscope TIFFs into channels
    \item \textbf{MATLAB detection} --- fit Gaussians to each spot (external code)
    \item \textbf{Analyze MATLAB output} --- filter puncta, export amplitudes
    \item \textbf{Calibrate \& plot} --- DLS calibration, then make figures
\end{enumerate}

% ═══════════════════════════════════════════════════════════════════
\section{Quick Primer: Using the Terminal}

Every step in this protocol involves typing commands into a terminal
(also called a \emph{command line}). Here's how to open one:

\begin{itemize}[leftmargin=*]
    \item \textbf{macOS:} Open Spotlight (\code{Cmd + Space}), type
    \code{Terminal}, press Enter.
    \item \textbf{Windows:} Press the Windows key, type
    \code{PowerShell}, click \textbf{Windows PowerShell}.
\end{itemize}

You'll see a window with a blinking cursor. This is where you type (or
paste) all the commands in this protocol. A few things to know:

\begin{itemize}[leftmargin=*]
    \item \code{cd} means ``\emph{change directory}.'' It's how you
    navigate to a folder. For example, \code{cd Desktop} moves you
    into your Desktop folder.
    \item You can paste commands with \code{Cmd+V} (Mac) or \code{Ctrl+V}
    / right-click (Windows). You don't have to type them out.
    \item Press \code{Enter} after each command to run it.
    \item If something goes wrong, read the error message --- it usually
    tells you what's missing. The Troubleshooting section at the end
    covers common errors.
\end{itemize}

\begin{tipbox}
\textbf{Tip.} If a command spans multiple lines with a backslash
\code{\textbackslash} (Mac) or backtick \code{`} (Windows) at the end
of each line, that's just one command split across lines for
readability. Paste the whole block at once --- the shell understands
the continuation.
\end{tipbox}

% ═══════════════════════════════════════════════════════════════════
\section{Prerequisites}

Before you start, you need these installed on your computer.

\subsection{Python}

Check if you already have Python by opening a terminal and typing:

\begin{lstlisting}
python --version
\end{lstlisting}

If you see something like \code{Python 3.11.5}, you're good. If you
get an error or see Python 2.x, install Python 3:

\begin{itemize}[leftmargin=*]
    \item \textbf{macOS:} Go to
    \href{https://www.python.org/downloads/}{python.org/downloads}
    and download the latest Python 3. Run the installer.
    \item \textbf{Windows:} Same website, same download.
    \textbf{Important:} during installation, check the box that says
    ``\emph{Add python.exe to PATH}'' --- this is on the very first
    screen of the installer.
\end{itemize}

\begin{notebox}
\textbf{Note.} On some Macs you may need to type \code{python3} instead
of \code{python}. If \code{python} doesn't work, try \code{python3}. If
\code{python3} works, use \code{python3} everywhere this protocol says
\code{python}.
\end{notebox}

\subsection{uv (Package Manager)}

\code{uv} is a fast Python package manager that replaces \code{pip}.
It keeps dependencies organized and avoids version conflicts.

\textbf{macOS / Linux:}

\begin{lstlisting}
curl -LsSf https://astral.sh/uv/install.sh | sh
\end{lstlisting}

\textbf{Windows PowerShell:}

\begin{lstlisting}
powershell -ExecutionPolicy ByPass -c "irm https://astral.sh/uv/install.ps1 | iex"
\end{lstlisting}

After installing, \textbf{close and reopen your terminal} so it picks
up the new command. Then verify it works:

\begin{lstlisting}
uv --version
\end{lstlisting}

You should see a version number like \code{uv 0.6.x}. If you get
``command not found,'' close and reopen the terminal again.

\subsection{MATLAB and CMEanalysis}

You need MATLAB installed with the CMEanalysis code from the Danuser
Lab. If MATLAB is already installed on your computer, skip to step 3.

\begin{enumerate}[leftmargin=*]
    \item \textbf{Download MATLAB.} Look for your institution's
    instructions. The software should be free for students and
    faculty.

    \textbf{Important:} when downloading, you will be prompted to
    select various toolboxes. These can be installed later, but some
    are needed for CMEanalysis. Select all of the following:

    \begin{itemize}
        \item Curve Fitting Toolbox
        \item Image Processing Toolbox
        \item Optimization Toolbox
        \item Statistics Toolbox
        \item Parallel Computing Toolbox
    \end{itemize}

    \item \textbf{Install missing toolboxes (if MATLAB is already
    installed).} Use the Add-On Explorer in the \textbf{Add-Ons} menu
    on the Home tab. You will need to restart MATLAB after each
    toolbox installs.

    \item \textbf{Install CMEanalysis.} Go to
    \href{https://github.com/DanuserLab/cmeAnalysis}{github.com/DanuserLab/cmeAnalysis},
    click the green \textbf{Code} button, and download the ZIP. Unzip
    the folder somewhere stable on your computer (e.g.\
    \code{Documents/cmeAnalysis-master}). \textbf{Remember the path
    where you put it} --- you will use it in Step 2.
\end{enumerate}

\subsection{DLS data}

You need a DLS measurement of your liposome stock from a Malvern
Zetasizer. Export the size distribution into an Excel spreadsheet.

\begin{notebox}
\textbf{Important: DLS data formatting.} Your Excel spreadsheet must
match the layout of the example file in the \code{test\_data/}
folder of this repository. Open
\code{test\_data/dls/example\_dls.xlsx} in Excel before formatting
your own --- copy the exact column layout. Specifically, the first
column must contain the section header \code{X Number} followed by
diameter bins, and subsequent columns hold the percentage data for
each measurement record.

The script \textbf{only uses the number distribution} (\code{X Number}
section). Other section headers like \code{X Intensity} and
\code{X Volume} are ignored, but they should still be present so the
file matches the standard Zetasizer export format.
\end{notebox}

% ═══════════════════════════════════════════════════════════════════
\section{Download the Code and Set Up}

Do these steps once. After this you only need to activate the virtual
environment (see \S\ref{sec:activate}) each time you open a new
terminal.

\subsection{Get the repository}

\textbf{Option A: with git} (if installed):

\begin{lstlisting}
git clone https://github.com/slele225/liposome-curvature-assay.git
\end{lstlisting}

\textbf{Option B: download a ZIP.} Go to
\href{https://github.com/slele225/liposome-curvature-assay}{github.com/slele225/liposome-curvature-assay},
click the green \textbf{Code} button, then \textbf{Download ZIP}.
Double-click to unzip. You should now have a folder called
\code{liposome-curvature-assay}.

\textbf{Move it somewhere you'll remember.} We recommend putting it on
your Desktop. The rest of this protocol assumes that location.

\subsection{Open a terminal in the repo folder}

\textbf{macOS / Linux:}

\begin{lstlisting}
cd ~/Desktop/liposome-curvature-assay
\end{lstlisting}

\textbf{Windows PowerShell:}

\begin{lstlisting}
cd $HOME\Desktop\liposome-curvature-assay
\end{lstlisting}

\begin{tipbox}
\textbf{Tip.} You can also type \code{cd} followed by a space, then
drag and drop the folder from Finder/File Explorer into the terminal
window. The full path will auto-fill.
\end{tipbox}

\subsection{Create a virtual environment}

This creates a private space for the Python packages so they don't
interfere with anything else on your computer.

\begin{lstlisting}
uv venv
\end{lstlisting}

You should see a message like \code{Creating virtual environment at:
.venv}. A new hidden folder called \code{.venv} now exists inside the
repo.

\subsection{Activate the virtual environment}\label{sec:activate}

\textbf{macOS / Linux:}

\begin{lstlisting}
source .venv/bin/activate
\end{lstlisting}

\textbf{Windows PowerShell:}

\begin{lstlisting}
.venv\Scripts\Activate
\end{lstlisting}

\textbf{How to tell it worked:} your terminal prompt should now start
with \code{(.venv)}. For example:

\begin{lstlisting}
(.venv) yourname@computer ~ %
\end{lstlisting}

\begin{warnbox}
\textbf{Important.} You need to activate the venv \textbf{every time}
you open a new terminal window. If you close the terminal and come
back later, repeat the \code{cd} step \emph{and} the activation step
above.
\end{warnbox}

\subsection{Install the Python packages}

\begin{lstlisting}
uv pip install -r requirements.txt
\end{lstlisting}

This downloads numpy, matplotlib, tifffile, h5py, openpyxl, pandas,
and scipy. It takes about a minute the first time, much faster on
later installs.

\subsection{Verify everything works}

\begin{lstlisting}
python plot_curvature.py --help
\end{lstlisting}

You should see a help message listing all the options for that script.
If you get an error, check that \code{(.venv)} appears at the start
of your terminal prompt. If not, re-activate the venv (\S\ref{sec:activate}).

% ═══════════════════════════════════════════════════════════════════
\section{Folder Structure: Where to Put Your Files}

The repo has two main folders for organizing your work. All your data
and outputs go inside the repo folder. Here is the structure (only
the most important files are shown):

\begin{lstlisting}
liposome-curvature-assay/
|-- data/                       <-- YOUR DATA GOES HERE
|   |-- dls_data.xlsx           <-- DLS Excel export
|   `-- experiment_name/        <-- one subfolder per experiment
|       `-- 488nm_5.0pct_580V_561nm_3.2pct_500V/
|           |-- cell1/
|           |   |-- ch1/        <-- lipid channel
|           |   `-- ch2/        <-- protein channel
|           `-- filtered_puncta_A_values.txt
|-- figures/                    <-- PLOTS SAVED HERE
|-- test_data/                  <-- example DLS file (reference format)
|-- prepare_input.py
|-- analyze_matlab.py
|-- dls_calibration.py
|-- plot_curvature.py
|-- plot_histograms.py
|-- plot_overlay.py
|-- plot_dls.py
|-- plot_dls_comparison.py
|-- requirements.txt
`-- README.md
\end{lstlisting}

\textbf{\code{data/}} --- Put your raw TIFF images in a subfolder here
(e.g.\ \code{data/march\_3\_experiment/}). The
\code{prepare\_input.py} script will create the MATLAB-ready folder
structure inside \code{data/}. Also put your DLS \code{.xlsx} file
here.

\textbf{\code{figures/}} --- This is where your final plots get
saved. The plotting scripts create this folder automatically when you
use \code{--save-dir figures/}.

\textbf{\code{test\_data/}} --- Reference data, including an example
DLS file showing the required spreadsheet format.

\begin{notebox}
\textbf{Note.} The \code{data/} and \code{figures/} folders are
gitignored, meaning their contents will not be uploaded if you use
git. Your microscopy data stays on your computer only.
\end{notebox}

% ═══════════════════════════════════════════════════════════════════
\section{Step 1: Prepare Images}

\textbf{What this does.} Takes your raw multi-frame TIFF files from
the microscope, splits them into separate per-channel folders, and
organizes them into the structure that MATLAB expects. Files are
automatically grouped by laser configuration. Olympus FV3000 metadata
is parsed for active lasers (any laser with transmissivity $> 0$),
each channel's excitation laser, and each channel's PMT voltage. The
output folder name encodes all three, e.g.\
\code{488nm\_5.0pct\_580V\_561nm\_3.2pct\_500V/}. Lambda-phase
channels are skipped. Channel folders \code{ch1}, \code{ch2}, etc.\
follow wavelength-ascending order (so \code{ch1} is always the lowest
active wavelength).

\subsection*{Before you start}

\begin{itemize}[leftmargin=*]
    \item Put all your raw \code{.tif} files from one experiment into
    a subfolder inside \code{data/}. For example:
    \code{data/march\_3\_experiment/}.
\end{itemize}

\subsection*{Run a dry-run first}

It's strongly recommended to first run with \code{--dry-run} to see
what folder structure the script will create, without actually
writing any files. Use this to verify that the metadata parsing matches
what you actually acquired.

\textbf{macOS / Linux:}

\begin{lstlisting}
python prepare_input.py \
    --input data/march_3_experiment \
    --output data/march_3_experiment_matlab \
    --crop 1 \
    --dry-run
\end{lstlisting}

\textbf{Windows PowerShell:}

\begin{lstlisting}
python prepare_input.py `
    --input data\march_3_experiment `
    --output data\march_3_experiment_matlab `
    --crop 1 `
    --dry-run
\end{lstlisting}

The script prints, for each TIFF, the active lasers and their
voltages, and the planned folder name. Read this carefully. If
anything looks wrong, fix it before running for real (most often, the
issue is metadata interpretation rather than the script).

\subsection*{Run for real}

If the dry-run output looks correct, repeat the same command without
\code{--dry-run}:

\textbf{macOS / Linux:}

\begin{lstlisting}
python prepare_input.py \
    --input data/march_3_experiment \
    --output data/march_3_experiment_matlab \
    --crop 1
\end{lstlisting}

\textbf{Windows PowerShell:}

\begin{lstlisting}
python prepare_input.py `
    --input data\march_3_experiment `
    --output data\march_3_experiment_matlab `
    --crop 1
\end{lstlisting}

\subsection*{All flags}

\begin{tabularx}{\linewidth}{l X}
\toprule
\textbf{Flag} & \textbf{Meaning} \\
\midrule
\flag{--input}    & Folder containing raw \code{.tif} files. \\
\flag{--output}   & Where to save the split channels (inside \code{data/}). \\
\flag{--frames}   & Comma-separated frame indices, one per active channel
                    in wavelength-ascending order (e.g.\ \code{0,1,2}).
                    Optional --- defaults to \code{0,1,...,N-1} per
                    TIFF based on its active channel count. \\
\flag{--crop}     & Center crop divisor. \code{1} = no crop, \code{2} =
                    center quarter. Use crop > 1 if your image has
                    dark borders or you want to focus on the central
                    region. \\
\flag{--dry-run}  & Parse metadata and print what would be created,
                    without writing files. \\
\bottomrule
\end{tabularx}

\subsection*{Expected output}

A new folder \code{data/march\_3\_experiment\_matlab/} containing
subfolders for each laser configuration, cell, and channel. For
example:

\begin{lstlisting}
data/march_3_experiment_matlab/
  `-- 488nm_5.0pct_580V_561nm_3.2pct_500V/
      |-- cell1/
      |   |-- ch1/      <-- lipid (488 nm)
      |   `-- ch2/      <-- protein (561 nm)
      `-- cell2/
          |-- ch1/
          `-- ch2/
\end{lstlisting}

% ═══════════════════════════════════════════════════════════════════
\section{Step 2: Run MATLAB Detection}

\textbf{What this does.} Fits 2D Gaussians to every sub-diffraction-
limit fluorescent spot in each channel image. This identifies
individual liposomes and measures their amplitudes (brightness),
backgrounds, and a quality flag.

These instructions use \textbf{CMEanalysis}, the external MATLAB code
base from the Danuser Lab that you installed in Prerequisites.

\subsection*{Procedure}

\begin{enumerate}[leftmargin=*]
    \item \textbf{Open MATLAB and navigate to the CMEanalysis folder.}
    In the MATLAB command window, type the following two commands.
    \textbf{Replace the path} with the actual location where you
    unzipped CMEanalysis (whatever path you remembered from
    Prerequisites). The path below is just an example.

    \begin{lstlisting}
cd('/path/to/cmeAnalysis-master/software')
addpath(genpath(pwd))
    \end{lstlisting}

    \item \textbf{Load the image folder.} In the MATLAB command
    window, run:

    \begin{lstlisting}
data = loadConditionData();
    \end{lstlisting}

    A file selection window will open. Select the \emph{condition
    folder} produced in Step 1 (the one named after the laser
    configuration, e.g.\
    \code{488nm\_5.0pct\_580V\_561nm\_3.2pct\_500V}).

    \item \textbf{Select channels.} The program will ask you to select
    the \emph{first (master) channel}. Select the lipid channel
    folder (this should be \code{ch1}). It will then ask for the
    remaining channels --- select \code{ch2} (and \code{ch3}, etc.\
    if you have more).

    \item \textbf{Enter fluorophore names.} The program will ask for a
    fluorescent marker for each channel. Type the fluorophore name
    matching that channel. Common ones include \code{texasred} for
    Texas Red and \code{egfp} for EGFP. If the program does not
    recognize the name, it will ask you to enter the emission
    wavelength in nm instead.

    \item \textbf{Run the detection.} The default mode determines the
    Gaussian-approximated PSF standard deviation directly from your
    images. Type:

    \begin{lstlisting}
runDetection(data);
    \end{lstlisting}

    If you have already run a detection on this folder and want to
    redo it (e.g.\ after changing parameters):

    \begin{lstlisting}
runDetection(data, 'Overwrite', true);
    \end{lstlisting}

    If you know what sigma value to use from separate experiments or
    calculations:

    \begin{lstlisting}
runDetection(data, 'Overwrite', true, 'Sigma', sigma_value);
    \end{lstlisting}

    Replace \code{sigma\_value} with your number. Note that sigma is
    the standard deviation, not the variance.

    \item \textbf{Wait for completion.} The code will show progress
    and take a few minutes. It reports the detected standard
    deviations for each channel across all images and shows graphs
    with this data --- you can close them.
\end{enumerate}

\subsection*{Verifying success}

When MATLAB finishes, verify that the master channel (\code{ch1})
now has a \code{Detection} subfolder. Note that \textbf{only the
master/lipid channel gets a Detection folder} --- \code{ch2} will
just contain the TIFF file. This is expected.

\begin{lstlisting}
data/march_3_experiment_matlab/488nm_5.0pct_580V_561nm_3.2pct_500V/
  `-- cell1/
      |-- ch1/
      |   `-- Detection/
      |       `-- detection_v2.mat   <-- this file must exist
      `-- ch2/
          `-- (TIFF image only, no Detection folder)
\end{lstlisting}

\begin{notebox}
\textbf{Note.} The \code{detection\_v2.mat} files contain a struct
called \code{frameInfo} with fields \code{A} (amplitude), \code{c}
(background), and \code{hval\_Ar} (hypothesis test value). The next
Python step reads these fields.
\end{notebox}

% ═══════════════════════════════════════════════════════════════════
\section{Step 3: Analyze MATLAB Output}

\textbf{What this does.} Reads the \code{.mat} detection files,
filters out noise and poorly-fit spots, and exports a single
tab-separated text file with the kept amplitude values for every
liposome that passed quality control.

\subsection*{Run the script}

\textbf{macOS / Linux (two-channel: lipid + protein):}

\begin{lstlisting}
python analyze_matlab.py \
    --input data/march_3_experiment_matlab/488nm_5.0pct_580V_561nm_3.2pct_500V \
    --channels ch1,ch2 \
    --lipid-channel ch1 \
    --k-std 2.0 \
    --output-name filtered_puncta_A_values.txt
\end{lstlisting}

\textbf{Windows PowerShell:}

\begin{lstlisting}
python analyze_matlab.py `
    --input data\march_3_experiment_matlab\488nm_5.0pct_580V_561nm_3.2pct_500V `
    --channels ch1,ch2 `
    --lipid-channel ch1 `
    --k-std 2.0 `
    --output-name filtered_puncta_A_values.txt
\end{lstlisting}

\subsection*{Lipid-only experiments}

If you only have a lipid channel (no protein), use \code{--channels
ch1} and the same channel for \code{--lipid-channel}:

\begin{lstlisting}
python analyze_matlab.py \
    --input data/march_3_experiment_matlab/488nm_5.0pct_580V_561nm_3.2pct_500V \
    --channels ch1 \
    --lipid-channel ch1
\end{lstlisting}

\subsection*{All flags}

\begin{tabularx}{\linewidth}{l X}
\toprule
\textbf{Flag} & \textbf{Meaning} \\
\midrule
\flag{--input}         & Condition folder (the laser-configuration
                         folder from Step 1). \\
\flag{--channels}      & Channel names matching folder names, comma-
                         separated. Use \code{ch1,ch2} for lipid +
                         protein. Use \code{ch1} for lipid-only. \\
\flag{--lipid-channel} & Which channel is the lipid (used to compute
                         the intensity threshold). Almost always
                         \code{ch1}. \\
\flag{--k-std}         & Threshold strictness. Keeps puncta where
                         amplitude $> \mu_c + k \cdot \sigma_c$, with
                         $\mu_c$ and $\sigma_c$ the mean and standard
                         deviation of the background. Default
                         \code{2.0}. Larger = stricter, fewer puncta. \\
\flag{--output-name}   & Output filename. Saved inside the
                         \code{--input} folder. Default
                         \code{filtered\_puncta\_A\_values.txt}. \\
\bottomrule
\end{tabularx}

\subsection*{Expected output}

The script prints a per-cell summary showing how many puncta were
kept vs.\ how many were detected. At the end, it saves a tab-
separated text file:

\begin{lstlisting}
data/march_3_experiment_matlab/488nm_5.0pct_580V_561nm_3.2pct_500V/
  filtered_puncta_A_values.txt
\end{lstlisting}

The file has columns: \code{source\_image}, \code{A\_ch1} (lipid
amplitude), \code{A\_ch2} (protein amplitude). Each row is one
liposome. Open it in Excel to spot-check.

% ═══════════════════════════════════════════════════════════════════
\section{Step 4: DLS Calibration}\label{sec:dls}

\textbf{What this does.} Reads your DLS Zetasizer spreadsheet and
your filtered puncta file from Step 3, then computes a single
\emph{conversion factor} that maps the square root of fluorescence
amplitude to physical diameter in nanometers. The script does this by
overlaying the fluorescence $\sqrt{A}$ distribution onto the DLS
number-weighted size distribution and finding the scalar that makes
them best match.

This is the standard SLiC calibration approach (Kunding 2008,
Hatzakis 2009, Bhatia 2009, Zeno 2018, Johnson 2025). A simpler
``ratio of means'' factor is also reported as a sanity check; the
overlay is more robust because it uses the full distribution shape.

\subsection*{Before you start}

Confirm your DLS spreadsheet matches the example format in
\code{test\_data/dls/}. See the note in the Prerequisites section
above.

\subsection*{Run the script}

\textbf{macOS / Linux:}

\begin{lstlisting}
python dls_calibration.py \
    --dls-input data/dls_data.xlsx \
    --fluor-input data/march_3_experiment_matlab/488nm_5.0pct_580V_561nm_3.2pct_500V/filtered_puncta_A_values.txt \
    --save-dir figures/
\end{lstlisting}

\textbf{Windows PowerShell:}

\begin{lstlisting}
python dls_calibration.py `
    --dls-input data\dls_data.xlsx `
    --fluor-input data\march_3_experiment_matlab\488nm_5.0pct_580V_561nm_3.2pct_500V\filtered_puncta_A_values.txt `
    --save-dir figures\
\end{lstlisting}

\subsection*{Optional: bootstrap variance estimate}

To estimate uncertainty in the conversion factor, add
\code{--bootstrap N}. The script resamples the puncta with
replacement N times and reports the coefficient of variation.

\begin{lstlisting}
python dls_calibration.py \
    --dls-input data/dls_data.xlsx \
    --fluor-input data/.../filtered_puncta_A_values.txt \
    --bootstrap 200 \
    --save-dir figures/
\end{lstlisting}

\subsection*{All flags}

\begin{tabularx}{\linewidth}{l X}
\toprule
\textbf{Flag} & \textbf{Meaning} \\
\midrule
\flag{--dls-input}    & Path to the Zetasizer \code{.xlsx} export. \\
\flag{--fluor-input}  & Path to \code{filtered\_puncta\_A\_values.txt}
                        from Step 3. \\
\flag{--bootstrap N}  & Bootstrap resamples for variance estimate.
                        Default \code{0} (off). Try \code{200}--\code{500}
                        to estimate uncertainty. \\
\flag{--bootstrap-k K}& Number of puncta to sample per bootstrap
                        iteration. Default = same as total puncta. \\
\flag{--lipid-col}    & Column name for lipid amplitude. Default
                        \code{A\_ch1}. \\
\flag{--save-dir}     & Save the overlay plots here (optional). \\
\bottomrule
\end{tabularx}

\subsection*{Output: what to copy down}

The script prints two key numbers near the bottom of its output. You
will use one of them in Step 5:

\begin{itemize}[leftmargin=*]
    \item \code{python plot\_curvature.py --conversion-factor X.XXXXXX
    ...} --- copy this number for use as
    \code{--conversion-factor}.
    \item Equivalently, an implied mean diameter for use as
    \code{--dls-mean-diameter}. Either choice gives the same plot.
\end{itemize}

It also saves a figure \code{figures/dls\_overlay\_calibration.png}
showing the DLS distribution (blue) and your fluorescence
distribution (red, dashed) on the same axis. They should overlap. If
they don't, your detection threshold (\code{--k-std} in Step 3) may
need adjustment.

% ═══════════════════════════════════════════════════════════════════
\section{Step 5: Plot Curvature Sorting}

\textbf{What this does.} Converts lipid amplitudes into physical
liposome diameters using the calibration from Step 4, computes
protein surface density for each liposome, and generates the final
scatter plot of \textbf{protein density vs.\ liposome diameter}. This
is the headline figure of the assay.

\subsection*{Run the script}

\textbf{macOS / Linux:}

\begin{lstlisting}
python plot_curvature.py \
    --input data/march_3_experiment_matlab/488nm_5.0pct_580V_561nm_3.2pct_500V/filtered_puncta_A_values.txt \
    --conversion-factor 1.234567 \
    --save-dir figures/
\end{lstlisting}

\textbf{Windows PowerShell:}

\begin{lstlisting}
python plot_curvature.py `
    --input data\march_3_experiment_matlab\488nm_5.0pct_580V_561nm_3.2pct_500V\filtered_puncta_A_values.txt `
    --conversion-factor 1.234567 `
    --save-dir figures\
\end{lstlisting}

Replace \code{1.234567} with the conversion factor printed by Step 4.

\subsection*{Alternative: use a mean diameter directly}

If you didn't run \code{dls\_calibration.py}, you can use a DLS mean
diameter directly (the script computes the ratio-of-means factor
internally):

\begin{lstlisting}
python plot_curvature.py \
    --input data/.../filtered_puncta_A_values.txt \
    --dls-mean-diameter 80.12 \
    --save-dir figures/
\end{lstlisting}

\subsection*{All flags}

\begin{tabularx}{\linewidth}{l X}
\toprule
\textbf{Flag} & \textbf{Meaning} \\
\midrule
\flag{--input}              & Filtered puncta file from Step 3. Accepts
                              multiple files (one plot per file). \\
\flag{--conversion-factor}  & Maps $\sqrt{A}$ to diameter in nm. From
                              \code{dls\_calibration.py}. Provide
                              this \emph{or} \code{--dls-mean-diameter},
                              not both. \\
\flag{--dls-mean-diameter}  & Alternative: mean DLS diameter in nm.
                              Internally uses ratio-of-means. \\
\flag{--lipid-col}          & Column for lipid amplitude. Default
                              \code{A\_ch1}. \\
\flag{--protein-col}        & Column for protein amplitude. Default
                              \code{A\_ch2}. \\
\flag{--bin-width}          & Diameter bin width in nm for the averaged
                              curve. Default \code{0.5}. Larger =
                              smoother. \\
\flag{--diameter-cutoff}    & Exclude puncta with diameter above this
                              value (nm). Useful for filtering sparse
                              large-diameter tail (optional). \\
\flag{--y-pad}              & Y-axis padding factor around bin means.
                              Default \code{0.3} (30\%). Lower values
                              like \code{0.1} zoom in tighter on the
                              trend. \\
\flag{--save-dir}           & Output directory for figures. \\
\bottomrule
\end{tabularx}

\subsection*{Expected output}

A PNG saved in \code{figures/}:
\code{488nm\_5.0pct\_580V\_561nm\_3.2pct\_500V\_\_protein\_density\_vs\_diameter.png}.

\subsection*{Interpreting the plot}

\begin{itemize}[leftmargin=*]
    \item \textbf{X-axis:} liposome diameter (nm). Smaller = higher
    curvature.
    \item \textbf{Y-axis:} protein surface density (amplitude per
    nm$^2$). Higher = more protein per unit membrane area.
    \item Light dots are individual liposomes; larger dots are
    bin-averaged means (the trend line you care about).
\end{itemize}

What the trend tells you:

\begin{itemize}[leftmargin=*]
    \item Bin averages going \textbf{up} toward smaller diameters
    $\Rightarrow$ your protein preferentially binds high-curvature
    membranes (positive curvature sensing).
    \item Flat line $\Rightarrow$ no curvature preference.
    \item Bin averages going \textbf{down} toward smaller diameters
    $\Rightarrow$ negative curvature sensing.
\end{itemize}

% ═══════════════════════════════════════════════════════════════════
\section{Step 6 (Optional): Plot Histograms}

A useful sanity-check tool. Plots amplitude histograms for each
channel and an estimated diameter distribution. Works for both
two-channel and lipid-only experiments. Most useful when you want to
verify the shape of your data before making the curvature plot.

\subsection*{Two-channel}

\begin{lstlisting}
python plot_histograms.py \
    --input data/.../filtered_puncta_A_values.txt \
    --lipid-col A_ch1 \
    --protein-col A_ch2 \
    --conversion-factor 1.234567 \
    --save-dir figures/
\end{lstlisting}

\subsection*{Lipid-only (omit \texttt{--protein-col})}

\begin{lstlisting}
python plot_histograms.py \
    --input data/.../filtered_puncta_A_values.txt \
    --lipid-col A_ch1 \
    --conversion-factor 1.234567 \
    --save-dir figures/
\end{lstlisting}

\subsection*{All flags}

\begin{tabularx}{\linewidth}{l X}
\toprule
\textbf{Flag} & \textbf{Meaning} \\
\midrule
\flag{--input}              & Filtered puncta file from Step 3. \\
\flag{--lipid-col}          & Lipid amplitude column. Default
                              \code{A\_ch1}. \\
\flag{--protein-col}        & Protein amplitude column. \emph{Omit
                              for lipid-only.} \\
\flag{--conversion-factor}  & From \code{dls\_calibration.py}. Omit to
                              skip the diameter plot. \\
\flag{--bins}               & Number of histogram bins. Default
                              \code{80}. \\
\flag{--transform}          & Amplitude transformation: \code{raw}
                              (default), \code{sqrt}, or \code{log\_sqrt}. \\
\flag{--save-dir}           & Output directory for figures. \\
\bottomrule
\end{tabularx}

% ═══════════════════════════════════════════════════════════════════
\section{Step 7 (Optional): Normalized Overlay Across Experiments}

If you have multiple experiments (different proteins, conditions,
replicates), overlay their normalized curvature-sorting curves on
one plot. By default each curve is normalized so its rightmost
(largest-diameter) bin equals 1, so the y-axis reads as
\emph{fold-enrichment at high curvature relative to flat membranes}
--- the same convention as Bhatia (2009) and Zeno (2018).

\subsection*{Run the script}

Each input is given as \code{path:conversion\_factor} (separated by
a colon). Run \code{dls\_calibration.py} \emph{separately for each
experiment} to get its own conversion factor.

\textbf{macOS / Linux:}

\begin{lstlisting}
python plot_overlay.py \
    --input data/cond1/filtered.txt:1.234 data/cond2/filtered.txt:1.567 \
    --labels "WT protein" "Mutant K58A" \
    --save-dir figures/
\end{lstlisting}

\textbf{Windows PowerShell:}

\begin{lstlisting}
python plot_overlay.py `
    --input data\cond1\filtered.txt:1.234 data\cond2\filtered.txt:1.567 `
    --labels "WT protein" "Mutant K58A" `
    --save-dir figures\
\end{lstlisting}

\subsection*{All flags}

\begin{tabularx}{\linewidth}{l X}
\toprule
\textbf{Flag} & \textbf{Meaning} \\
\midrule
\flag{--input}            & Files with conversion factors, as
                            \code{file.txt:factor} pairs. \\
\flag{--labels}           & Custom legend labels (default = parent
                            folder name of each file). \\
\flag{--lipid-col}        & Lipid column. Default \code{A\_ch1}. \\
\flag{--protein-col}      & Protein column. Default \code{A\_ch2}. \\
\flag{--bin-width}        & Diameter bin width in nm. Default \code{0.5}. \\
\flag{--normalize-to}     & How to normalize curves:
                            \code{rightmost} (default, matches
                            Bhatia/Zeno), \code{leftmost} (for
                            negative curvature sensors),
                            \code{minimum} (always $\geq 1$), or
                            \code{none}. \\
\flag{--diameter-cutoff}  & Exclude puncta above this diameter (nm). \\
\flag{--y-pad}            & Y-axis padding factor. Default \code{0.3}. \\
\flag{--output-name}      & Output filename. Default
                            \code{normalized\_curvature\_overlay.png}. \\
\flag{--save-dir}         & Output directory. \\
\bottomrule
\end{tabularx}

% ═══════════════════════════════════════════════════════════════════
\section{Step 8 (Optional): DLS Distribution Plot}

Quick standalone visualization of your DLS distribution.

\begin{lstlisting}
python plot_dls.py data/dls_data.xlsx
python plot_dls.py data/dls_data.xlsx --log
python plot_dls.py data/dls_data.xlsx --distribution intensity
python plot_dls.py data/dls_data.xlsx --zoom-pct 100
\end{lstlisting}

\begin{tabularx}{\linewidth}{l X}
\toprule
\textbf{Flag} & \textbf{Meaning} \\
\midrule
\flag{input}            & Path to Zetasizer \code{.xlsx} (positional). \\
\flag{--distribution}   & \code{number} (default) or \code{intensity}. \\
\flag{--log}            & Plot $\log(\text{diameter})$ instead of raw. \\
\flag{--zoom-pct}       & Percent of data to show (default \code{95}). \\
\bottomrule
\end{tabularx}

% ═══════════════════════════════════════════════════════════════════
\section{Step 9 (Optional): DLS vs Fluorescence Comparison}

Side-by-side comparison of the DLS distribution with the
fluorescence $\sqrt{A}$ distributions for one or more channels. One
panel per channel. Use this to verify that your detected liposome
population matches the DLS measurement.

\begin{lstlisting}
python plot_dls_comparison.py \
    --dls-input data/dls_data.xlsx \
    --fluor-input data/.../filtered_puncta_A_values.txt \
    --channels 0 Lipid 1 EGFP \
    --zoom-pct 95 \
    --bins 200 \
    --save-dir figures/
\end{lstlisting}

\begin{tabularx}{\linewidth}{l X}
\toprule
\textbf{Flag} & \textbf{Meaning} \\
\midrule
\flag{--dls-input}        & DLS \code{.xlsx} file. \\
\flag{--fluor-input}      & Filtered puncta file. \\
\flag{--channels}         & Index/label pairs: \code{0 Lipid 1 EGFP}.
                            Index 0 = \code{A\_ch1}, 1 = \code{A\_ch2}. \\
\flag{--dls-distribution} & \code{number} (default) or \code{intensity}
                            for the DLS panel. \\
\flag{--bins}             & Bins for the fluorescence histogram (default
                            \code{100}). \\
\flag{--zoom-pct}         & Percent of data to show (default \code{100}). \\
\flag{--save-dir}         & Output directory. \\
\bottomrule
\end{tabularx}

% ═══════════════════════════════════════════════════════════════════
\section{Quick Reference: Every-Time Checklist}

After the initial setup is done, here's what you do each time you
have new data:

\begin{enumerate}[leftmargin=*]
    \item Open a terminal.
    \item Navigate to the repo:

    \begin{itemize}
        \item Mac: \code{cd \textasciitilde/Desktop/liposome-curvature-assay}
        \item Windows: \code{cd \$HOME$\backslash$Desktop$\backslash$liposome-curvature-assay}
    \end{itemize}

    \item Activate the venv:

    \begin{itemize}
        \item Mac: \code{source .venv/bin/activate}
        \item Windows: \code{.venv$\backslash$Scripts$\backslash$Activate}
    \end{itemize}

    \item Copy your raw TIFFs into a new subfolder inside
    \code{data/}.
    \item Run \code{prepare\_input.py} (Step 1) --- with
    \code{--dry-run} first to verify.
    \item Run MATLAB detection (Step 2).
    \item Run \code{analyze\_matlab.py} (Step 3).
    \item Run \code{dls\_calibration.py} (Step 4) --- copy down the
    conversion factor.
    \item Run \code{plot\_curvature.py} (Step 5) with that factor and
    \code{--save-dir figures/}.
    \item Find your plot in the \code{figures/} folder.
\end{enumerate}

% ═══════════════════════════════════════════════════════════════════
\section{Troubleshooting}

This section covers errors you are likely to encounter. Most issues
trace back to one of three things: the virtual environment is not
active, a path is wrong, or the data file isn't where the script
expects.

\subsection{Setup errors}

\textbf{``\code{python} is not recognized as an internal or external
command'' (Windows).} Python wasn't added to \code{PATH} during
install. Either reinstall Python and check ``Add python.exe to
PATH,'' or use the full path to \code{python.exe}.

\textbf{``\code{python3} is not recognized'' (Windows).} On Windows,
use \code{python} (not \code{python3}). The \code{python3} convention
is Mac/Linux only.

\textbf{``\code{command not found: python}'' (Mac).} Try
\code{python3}. On some Macs, \code{python} isn't linked but
\code{python3} is. If neither works, install Python (see
Prerequisites).

\textbf{``\code{command not found: uv}.''} Close and reopen your
terminal after installing uv. If still not found, reinstall uv with
the commands in Prerequisites.

\textbf{``\code{No module named tifffile/numpy/pandas/...}.''} Your
virtual environment is not active. Run
\code{source .venv/bin/activate} (Mac) or
\code{.venv$\backslash$Scripts$\backslash$Activate} (Windows). Look
for \code{(.venv)} at the start of your prompt.

\textbf{Activation script can't be run on Windows (PowerShell
execution policy error).} Run this once, then try activating again:

\begin{lstlisting}
Set-ExecutionPolicy -ExecutionPolicy RemoteSigned -Scope CurrentUser
\end{lstlisting}

\textbf{``Permission denied'' on Mac.} For the uv install script:
prepend \code{sudo}. For \code{.venv}: delete the \code{.venv}
folder and recreate it with \code{uv venv}.

\subsection{File and path errors}

\textbf{``\code{FileNotFoundError: [...] filtered\_puncta\_A\_values.txt}.''}
Either the path is wrong, or Step 3 didn't write the file. Paste
your path into Finder/File Explorer to confirm it exists. If the
path contains spaces, wrap it in double quotes.

\textbf{``\code{No .tif files found}.''} \code{--input} points to
the wrong folder. It should be the folder containing your TIFFs
directly, not a parent or sibling.

\textbf{``\code{missing detection file}.''} CMEanalysis hasn't been
run on that cell, or it failed silently. Open the cell folder in
Finder/Explorer and check whether \code{ch1/Detection/} exists with
a \code{detection\_v2.mat} inside.

\textbf{Path with spaces causes an error.} Wrap the entire path in
double quotes:

\begin{lstlisting}
python prepare_input.py --input "data/march 3 experiment" ...
\end{lstlisting}

\subsection{Pipeline-stage errors}

\textbf{Step 1: ``\code{Frames out of range}.''} The frame indices
in \code{--frames} don't match what's in the TIFF. Open one TIFF in
ImageJ and check how many frames it has. Frame indices start at 0.

\textbf{Step 2: MATLAB detection produced no \code{Detection/}
folder.} Make sure you set \code{ch1} as the master channel in
\code{loadConditionData()}. Only the master channel gets a
\code{Detection/} folder.

\textbf{Step 3: ``No puncta passed the filter.''} The threshold
(\code{--k-std}) is too strict, or signal is too low. Try
\code{--k-std 1.5} or \code{--k-std 1.0}.

\textbf{Step 4: ``\code{Could not find 'X Number' section in
spreadsheet}.''} Your DLS Excel file is not in the expected format.
Open \code{test\_data/dls/example\_dls.xlsx} and copy that exact
layout. The first column must contain the section header
\code{X Number} followed by diameter bins.

\textbf{Step 4: implied mean diameter is way off (e.g.\ $<10$ nm or
$>1000$ nm).} The conversion likely failed because of a Step 3
issue. Run \code{plot\_histograms.py} to look at your fluorescence
distribution --- it should be roughly unimodal and positive. If it
looks bad, revisit your \code{--k-std} threshold in Step 3.

\textbf{Step 5: x-axis range doesn't match DLS.} You're applying the
conversion factor wrong, or the implied conversion factor is bad.
Make sure you copied the exact number printed by
\code{dls\_calibration.py} into \code{--conversion-factor}, not a
derived value.

\subsection{Plot quality issues}

\textbf{Plot looks flat / no trend visible.} The y-axis is being
auto-scaled to outliers. Add \code{--y-pad 0.1} to zoom in on the
trend, or use \code{--diameter-cutoff 120} to drop the noisy
large-diameter tail.

\textbf{Curve is very noisy.} Increase \code{--bin-width} (try
\code{1.0} or \code{2.0}). Larger bins average over more puncta and
smooth out noise.

\textbf{Plot's diameter range is unreasonable.} See ``Step 5: x-axis
range doesn't match DLS'' above.

\subsection{General}

\textbf{I closed my terminal and now nothing works.} You need to
\code{cd} back to the repo folder and activate the venv again. See
the Quick Reference checklist.

\textbf{``UnicodeDecodeError'' or weird text errors.} Your DLS
\code{.xlsx} may have been opened in Google Sheets and saved as
\code{.csv}. Re-export from the Malvern software as a fresh
\code{.xlsx}.

% ═══════════════════════════════════════════════════════════════════
\section{Common Pitfalls}

A few extra things to watch out for, gathered from real users running
this pipeline:

\begin{itemize}[leftmargin=*]
    \item \textbf{Forgetting to activate the venv.} The single most
    common error. Always check that your prompt starts with
    \code{(.venv)}.
    \item \textbf{Mixing up \code{--conversion-factor} and
    \code{--dls-mean-diameter}.} They are not the same number. The
    conversion factor is on the order of $1$ to $10$ (units of
    nm/$\sqrt{A}$); the mean diameter is in nanometers (typically
    50--200 nm). If your plot's x-axis is way off, you probably
    swapped them.
    \item \textbf{Re-using an old conversion factor for new data.}
    Each liposome stock has its own DLS distribution and its own
    calibration. Always re-run \code{dls\_calibration.py} for new
    data.
    \item \textbf{Comparing two experiments without separate
    calibrations.} For \code{plot\_overlay.py}, each input file
    needs its own conversion factor. Don't reuse one factor across
    multiple experiments.
    \item \textbf{Saving the plot to the wrong folder.} If you
    forget \code{--save-dir figures/}, the plot saves in your
    current working directory. Always pass the flag.
    \item \textbf{Spaces in folder names without quotes.} Wrap any
    path with spaces in double quotes:
    \code{"data/march 3 experiment"}.
    \item \textbf{Editing files in Excel and saving as .csv.} Some
    DLS workflows tempt you to save as CSV. Don't ---
    \code{dls\_calibration.py} expects \code{.xlsx} with the
    Zetasizer section headers in column A.
\end{itemize}

% ═══════════════════════════════════════════════════════════════════
\section{Further Reading}

\begin{itemize}[leftmargin=*]
    \item \code{README.md} --- terse developer reference for every
    command.
    \item \code{DLS\_Calibration\_Notes.md} --- theoretical
    background on the distribution-overlay calibration: why it
    works, how it relates to published SLiC protocols, what the
    reported statistics mean.
    \item Original assay references: Kunding et al.\ (2008) Biophys
    J 95:1176; Hatzakis et al.\ (2009) Nat Chem Biol 5:835; Bhatia
    et al.\ (2009) EMBO J 28:3303; Zeno et al.\ (2018) Nat Commun
    9:4152; Johnson et al.\ (2025) Commun Biol 8:1179.
\end{itemize}

\end{document}